\documentclass{article}
\usepackage[dvipsnames]{xcolor}

\newcommand{\cmd}[1]{\texttt{\textcolor{blue}{#1}}}

\begin{document}

\centerline{\Large\textbf{Cheatsheet}}

\vspace{1cm}

\noindent Commandes pour toy\_example.py : \\
\begin{verbatim}
toy_example.py [-h] -p PREFIX [-s STATES] [-l {bin,pos,state,ml}] [-d] [-m] [-a]
[-i {brute,pred,voteF,voteQ,find}] [-dot] [-pm PATH_MODEL] [-dp PATH_DATA]    
\end{verbatim}
Pour avoir de l'aide sur les arguments, la commande \cmd{toy\_example.py -h} affiche la description de chaque argument.
Renseigner un prefixe inconnu pour l'argument \cmd{-p} créera un nouveau dataset à partir d'un automate construit 
aléatoirement, et sauvegardera le dataset et l'automate dans le dossier \texttt{test/} par défaut.
Cependant, il est possible de construire un automate en particulier grâce à build\_auto.py :
\begin{verbatim}
    build_auto.py [-h] [-a ALPHABET] [-m MOTS] [-n NOM] [-t]
\end{verbatim}

\subsection*{Exemples d'utilisation}

Pour toy\_example.py, la commande suivante :
{\color{blue}
\begin{verbatim}
    toy_example.py -p FSM4 -s 100 -l ml -i pred -d -m -a
\end{verbatim}}
\begin{itemize}
    \item créée un dataset à partir de l'automate FSM4
    \item entraine un modèle de GRU en "multi-label"
    \item extraie un automate de 100 états à partir du modèle entrainé
\end{itemize}
Tandis que la commande :
{\color{blue}
\begin{verbatim}
    toy_example.py -p FSM4 -s 100 -l ml -dot
\end{verbatim}}
\begin{itemize}
    \item test le modèle de GRU en "multi-label" sur l'automate FSM4
    \item test l'automate extrait à partir du modèle de GRU
    \item sauvegarde l'automate au format .dot
\end{itemize}
{\color{red}Attention :} lève une exception si le modèle ou l'automate ne sont pas créés.

\vspace{0.3cm}
\noindent Pour créer un automate qui reconnait abc,abd avec build\_auto.py :
{\color{blue}
\begin{verbatim}
    build_auto.py -a abcd -m abc,abd -n auto_custom
\end{verbatim}}
\noindent Il est possible de visualiser l'automate sans le sauvegarder grâce à \cmd{-t}.




\newpage
\subsection*{Hyperparamètres}
Les hyperparamètres gravient autour des différentes méthodes pour
étiqueter les données ou pour trouver l'état initial de l'automate extrait.
\\Méthodes d'étiquetage (argument \cmd{-l}) :
\begin{itemize}
    \item \textit{bin} : (binaire), 1 si la séquence est acceptée par l'automate, 0 sinon.
    \item \textit{pos} : (multi-classe), position de la lettre dans la séquence, 0 si la lettre n'appartient à une sequence.
    \item \textit{state} : (multi-classe), nom de l'état de l'automate atteint après la lecture de la lettre.
    \item \textit{ml} : (multi-label), liste des positions de la lettre pour chaque séquence.
\end{itemize}

\noindent Méthodes pour trouver l'état initial de l'automate extrait (argument \cmd{-i}) :
\begin{itemize}
    \item \textit{brute} : (brute-force), on créer un cluster en plus à partir d'un vecteur nul pour représenter l'état initial de l'automate extrait.
    \item \textit{pred} : (prédiction), la prédiction de kmeans nous permets d'associer le vecteur nul à un cluster existant.
    \item \textit{voteF} : (vote finaux), on regarde les états successeurs des états finaux de l'automate extrait, celui qui est le plus représenté est considéré comme l'état initial.
    \item \textit{voteQ} : (vote états), même méthode mais pour tous les états de l'automate extrait.
    \item \textit{find} : (recherche), regarde le dernier état atteint par des séquences aléatoires, celui qui est le plus représenté est considéré comme l'état initial.
\end{itemize}
\end{document}